% ============================================================================
% sec06.tex — §6 Ideal & Real Gases, Kinetic Theory, Maxwell-Boltzmann Speeds,
%             Equipartition & Van der Waals (Target: 10 pages)
% ============================================================================
\binderchapter{6}{Ideal \& Real Gases, Kinetic Theory, Maxwell--Boltzmann Speeds, Equipartition \& Van der Waals}{Gases \& Kinetic Theory}{sec:06}
\scopebanner{REQUIRED}[2027 Rule 5.c.ii \& 5.c.iii (All Tournament Levels) \textbar\ Real-Gas \& Quantum Freeze-Out Extensions]

\idxtopic{ideal gas law}\idxtopic{kinetic theory of gases}\idxtopic{Maxwell--Boltzmann distribution}\idxtopic{equipartition theorem}\idxtopic{van der Waals equation}\idxtopic{Dalton's law of partial pressures}\idxtopic{Graham's law of effusion}

Gas-law items appear on nearly every exam in our 8-exam corpus (\sref{sec:22}) \src{S25,S26,S27,S28}, spanning foundational pressure/temperature unit traps through multi-compartment piston equilibrium and quantum heat-capacity freeze-out. Every governing relation in this chapter is presented in closed algebraic form under Rule~\texttt{5.d.iv} \src{S01,S06}.

% ============================================================================
\section{Macroscopic Gas Variables: Absolute vs.\ Gauge Pressure, Manometry \& STP Standards}\label{sec:06.1}
% ============================================================================
\idxtopic{pressure!absolute vs. gauge}\idxtopic{manometer}\idxtopic{barometer}\idxtopic{STP vs. NTP vs. SATP}

Gas equations of state respond strictly to \textbf{absolute pressure} $P_{\mathrm{abs}}$ and \textbf{absolute temperature} $T$ (in $\si{\kelvin}$ or $^\circ\mathrm{R}$). Most mechanical pressure gauges (Bourdon tubes, tire gauges, boiler dials) measure pressure relative to the local ambient atmosphere $P_{\mathrm{atm}}$ \src{S13,S25}:
\begin{equation}\label{eq:06:pressure-abs}
  P_{\mathrm{abs}} = P_{\mathrm{atm}} + P_{\mathrm{gauge}}, \qquad P_{\mathrm{abs}} = P_{\mathrm{atm}} - P_{\mathrm{vac}}
\end{equation}
\idxformula{Absolute pressure: $P_{\mathrm{abs}} = P_{\mathrm{atm}} + P_{\mathrm{gauge}}$}
where $P_{\mathrm{atm}} = \SI{101325}{\pascal} = \SI{101.325}{\kilo\pascal} = \SI{1.01325}{\bar} = \SI{760}{\text{Torr}} = \SI{14.6959}{\text{psi}}$ is the standard sea-level atmosphere \src{S07,S08}.

\begin{formulacard}[Hydrostatic Manometry, Barometry \& Reference Conditions \src{S07,S08,S13}]
\begin{itemize}
  \item \textbf{Hydrostatic Column Pressure Difference:} For an incompressible manometer fluid of density $\rho_f$ in gravitational field $g_n = \SI{9.80665}{\meter\per\square\second}$ with vertical height difference $h$:
  \begin{equation}\label{eq:06:manometer}
    \Delta P = \rho_f g_n h, \qquad P_{\mathrm{gas}} = \begin{cases} \rho_f g_n h & \text{(closed-end vacuum manometer, } P_{\mathrm{ref}}\approx 0\text{)}\\ P_{\mathrm{atm}} \pm \rho_f g_n h & \text{(open-end U-tube: $+$ if open limb is higher; $-$ if gas limb is higher)} \end{cases}
  \end{equation}
  \idxformula{Manometer pressure: $P_{\mathrm{gas}} = P_{\mathrm{atm}} \pm \rho_f g_n h$}
  \item \textbf{Fluid-Column Equivalence ($h_1\rho_1 = h_2\rho_2$):} Using $\rho_{\mathrm{Hg}}(\SI{0}{\celsius}) = \SI{13595.1}{\kilogram\per\cubic\meter}$ and $\rho_{\mathrm{H_2O}}(\SI{4}{\celsius}) = \SI{999.97}{\kilogram\per\cubic\meter}$, a column of mercury of height $h_{\mathrm{Hg}}$ corresponds to an equivalent water column $h_{\mathrm{H_2O}} = (\rho_{\mathrm{Hg}}/\rho_{\mathrm{H_2O}})h_{\mathrm{Hg}} = 13.595\,h_{\mathrm{Hg}}$ ($\SI{1}{\text{atm}} = \SI{760}{\milli\meter\text{Hg}} = \SI{10.332}{\meter\text{H}_2\text{O}}$).
\end{itemize}
\par\smallskip\noindent
\footnotesize
\begin{tabularx}{\linewidth}{@{}l l l l X@{}}
\toprule
\textbf{Standard Label} & \textbf{Temperature $T$} & \textbf{Pressure $P$} & \textbf{Ideal $V_m = RT/P$} & \textbf{Authority \& Competition Usage} \\
\midrule
\textbf{IUPAC STP (Current)} & $\SI{0}{\celsius} = \SI{273.15}{\kelvin}$ & $\SI{100}{\kilo\pascal} = \SI{1.000}{\bar}$ & $\SI{22.71095}{\liter\per\mole}$ & IUPAC since 1982; NIST CODATA 2022 \src{S07} \\
\textbf{Traditional STP} & $\SI{0}{\celsius} = \SI{273.15}{\kelvin}$ & $\SI{101.325}{\kilo\pascal} = \SI{1.000}{\text{atm}}$ & $\SI{22.41397}{\liter\per\mole}$ & Classic high-school textbooks \& invitational MCQs \src{S13,S29} \\
\textbf{SATP (Thermochem.)} & $\SI{25}{\celsius} = \SI{298.15}{\kelvin}$ & $\SI{100}{\kilo\pascal} = \SI{1.000}{\bar}$ & $\SI{24.78957}{\liter\per\mole}$ & Standard ambient tables ($\Delta_f H^\circ, S^\circ, \Delta_f G^\circ$) \src{S10,S16} \\
\textbf{NTP (Engineering)} & $\SI{20}{\celsius} = \SI{293.15}{\kelvin}$ & $\SI{101.325}{\kilo\pascal} = \SI{1.000}{\text{atm}}$ & $\SI{24.05485}{\liter\per\mole}$ & NIST / industrial compressed-gas metering \\
\bottomrule
\end{tabularx}
\end{formulacard}

\begin{workedexample}[Marcet Boiler Gauge-to-Absolute Specific Volume \& Compressibility Audit]{2027 Cornell E1 Q27c \textbar\ Tier 2}{ex:06.1}
In a Regnault/Marcet steam-boiler experiment at $T = \SI{120.0}{\celsius}$ ($\SI{393.15}{\kelvin}$), the boiler pressure gauge reads $P_{\mathrm{gauge}} = \SI{97.35}{\kilo\pascal}$ when the local barometric pressure is $P_{\mathrm{atm}} = \SI{101.325}{\kilo\pascal}$. The measured specific volume of saturated water vapor from Clapeyron slope analysis is $v_{g,\mathrm{meas}} = \SI{0.8919}{\cubic\meter\per\kilogram}$ \src{S09,S25}. (a) Calculate the ideal-gas specific volume $v_{g,\mathrm{ideal}}$ of steam ($M = \SI{18.0153}{\gram\per\mole}$) at this state. (b) Determine the fractional error incurred if a competitor forgets to add $P_{\mathrm{atm}}$. (c) Compute the real-gas compressibility factor $Z = P_{\mathrm{abs}}v_{g,\mathrm{meas}}/(R_s T)$ and explain why $Z < 1$.
\par\smallskip\noindent
\textbf{Solution.}
\begin{enumerate}
  \item \textbf{Absolute pressure and specific gas constant:}
  \[
    P_{\mathrm{abs}} = P_{\mathrm{atm}} + P_{\mathrm{gauge}} = 101.325 + 97.35 = \SI{198.675}{\kilo\pascal} = \SI{1.98675e5}{\pascal}, \qquad R_s = \frac{R}{M} = \frac{\SI{8.3144626}{\joule\per\mole\per\kelvin}}{\SI{0.0180153}{\kilogram\per\mole}} = \SI{461.522}{\joule\per\kilogram\per\kelvin}
  \]
  \item \textbf{(a) Ideal-gas specific volume at $\SI{393.15}{\kelvin}$:}
  \[
    v_{g,\mathrm{ideal}} = \frac{R_s T}{P_{\mathrm{abs}}} = \frac{461.522 \times 393.15}{1.98675\times 10^5} = \mathbf{\SI{0.91329}{\cubic\meter\per\kilogram}}
  \]
  \item \textbf{(b) Gauge-pressure trap error:} Using $P_{\mathrm{gauge}} = \SI{97.35}{\kilo\pascal}$ gives $v_{\mathrm{wrong}} = \frac{461.522\times 393.15}{9.735\times 10^4} = \SI{1.8639}{\cubic\meter\per\kilogram}$, an overestimate of $+104.1\%$ (more than double the true volume!).\idxerr{gauge vs. absolute pressure in ideal gas law}
  \item \textbf{(c) Compressibility factor $Z$:}
  \[
    Z = \frac{v_{g,\mathrm{meas}}}{v_{g,\mathrm{ideal}}} = \frac{0.8919}{0.91329} = \mathbf{0.9766} < 1
  \]
  At saturation near $\SI{2}{\bar}$, polar $\mathrm{H_2O}$ molecules experience strong attractive dipole and hydrogen-bonding forces ($a$ term), pulling molecules closer together than an ideal gas and reducing $v_g$ by $2.34\%$.
\end{enumerate}
\verifiedby{Direct evaluation from $P_{\mathrm{abs}}V_m = ZRT$ with $V_{m,\mathrm{meas}} = \SI{0.016068}{\cubic\meter\per\mole}$ gives $Z = (198675\times 0.016068)/(8.31446\times 393.15) = 0.9766$ (\texttt{check\_sec06\_09.py::ex06\_1}).}
\end{workedexample}

% ============================================================================
\section{The Ideal Gas Equation of State, Classical Limit Laws \& Multi-Bulb Mass Ledgers}\label{sec:06.2}
% ============================================================================
\idxtopic{Boyle's law}\idxtopic{Charles's law}\idxtopic{Gay-Lussac's law}\idxtopic{Avogadro's law}\idxtopic{gas constant $R$}

In the dilute limit ($P \ll P_c$ and $T \gtrsim 2T_b$), all gases obey the universal \textbf{Ideal Gas Equation of State} \src{S07,S13}:
\begin{equation}\label{eq:06:ideal-gas}
  PV = nRT = N k_B T = m R_s T, \qquad \rho = \frac{m}{V} = \frac{PM}{RT} = \frac{P}{R_s T}, \qquad n_{\#}=\frac{N}{V}=\frac{P}{k_B T}
\end{equation}
\idxformula{Ideal gas law: $PV = nRT = Nk_BT = mR_sT$}\idxformula{Ideal gas density: $\rho = PM/(RT)$}
where $n = m/M = N/N_A$ is the amount of substance ($\si{\mole}$), $M$ is molar mass ($\si{\kilogram\per\mole}$ in SI formulas!), $R = N_A k_B = \SI{8.314462618}{\joule\per\mole\per\kelvin}$, $R_s = R/M$ is the specific gas constant ($\si{\joule\per\kilogram\per\kelvin}$), and $n_{\#} = N/V$ is number density ($\si{\per\cubic\meter}$; at $1\text{ atm}, \SI{273.15}{\kelvin}$, Loschmidt's constant is $n_0 = \SI{2.68678e25}{\per\cubic\meter}$ \src{S07}).

\begin{twocol}[0.51]
\begin{formulacard}[Classical Gas Laws \& Valve/Leak Ledgers \src{S13}]
\par\smallskip\noindent
\footnotesize
\begin{tabularx}{\linewidth}{@{}l l X@{}}
\toprule
\textbf{Law / Invariant} & \textbf{Held Fixed} & \textbf{Algebraic Ratio Form} \\
\midrule
\textbf{Combined Gas Law} & — & $\dfrac{P_1 V_1}{n_1 T_1} = \dfrac{P_2 V_2}{n_2 T_2} = R$ \\[5pt]
\textbf{Boyle's Law} (1662) & $n, T$ & $P_1 V_1 = P_2 V_2$ (hyperbola) \\[2pt]
\textbf{Charles's Law} (1787) & $n, P$ & $V_1/T_1 = V_2/T_2$ ($\beta_V = 1/T$) \\[2pt]
\textbf{Amontons / Gay-Lussac} & $n, V$ & $P_1/T_1 = P_2/T_2$ \\[2pt]
\textbf{Avogadro's Law} (1811) & $P, T$ & $V_1/n_1 = V_2/n_2 = V_m$ \\[2pt]
\textbf{Density-State Ratio} & Gas $M$ & $\dfrac{P_1}{\rho_1 T_1} = \dfrac{P_2}{\rho_2 T_2} = R_s$ \\
\bottomrule
\end{tabularx}
\par\smallskip\noindent
\textbf{Connected Rigid Bulbs at Different Temperatures ($T_A, T_B$):} When a valve opens between bulbs of volumes $V_A$ (held at $T_A$) and $V_B$ (held at $T_B$), mechanical equilibrium enforces a \textbf{uniform final pressure $P_f$} across the capillary while total moles are conserved ($n_{\mathrm{tot}} = n_{A,i} + n_{B,i} = n_{A,f} + n_{B,f}$):
\begin{equation}\label{eq:06:two-bulbs}
  \frac{P_{A,i}V_A}{T_{A,i}} + \frac{P_{B,i}V_B}{T_{B,i}} = P_f\left(\frac{V_A}{T_{A,f}} + \frac{V_B}{T_{B,f}}\right)
\end{equation}
\idxformula{Connected gas bulbs: $P_f = n_{\mathrm{tot}}R/(V_A/T_{A,f} + V_B/T_{B,f})$}
\end{formulacard}
\nextcol[0.47]
\begin{shortcut}[7 Competition Forms of $R$ \src{S07,S08}]
Pick the value of $R$ matching the problem's pressure and volume units to avoid conversion errors:
\par\smallskip\noindent
\footnotesize
\begin{tabularx}{\linewidth}{@{}l X@{}}
\toprule
\textbf{Numerical Value of $R$} & \textbf{Exact Unit Combination} \\
\midrule
$\mathbf{8.314463}$ & $\si{\joule\per\mole\per\kelvin} = \si{\pascal\cubic\meter\per\mole\per\kelvin}$ \\
$\mathbf{8.314463}$ & $\si{\liter\kilo\pascal\per\mole\per\kelvin}$ ($1\,\si{\liter\kilo\pascal} = \SI{1}{\joule}$) \\
$\mathbf{0.0831446}$ & $\si{\liter\bar\per\mole\per\kelvin}$ ($1\,\si{\liter\bar} = \SI{100}{\joule}$) \\
$\mathbf{0.0820574}$ & $\si{\liter\text{atm}\per\mole\per\kelvin}$ ($1\,\si{\liter\text{atm}} = \SI{101.325}{\joule}$) \\
$\mathbf{62.36360}$ & $\si{\liter\text{Torr}\per\mole\per\kelvin}$ (or $\si{\liter\milli\meter\text{Hg}\per\mole\per\kelvin}$) \\
$\mathbf{1.987204}$ & $\si{\text{cal}_{\mathrm{th}}\per\mole\per\kelvin}$ ($\approx \SI{2.0}{\text{cal}\per\mole\per\kelvin}$) \\
$\mathbf{8.31446\times 10^7}$ & $\si{\text{erg}\per\mole\per\kelvin}$ ($1\,\si{\joule} = 10^7\,\si{\text{erg}}$) \\
\bottomrule
\end{tabularx}
\end{shortcut}
\begin{trapbox}[Open-Room Heating: Internal Energy Does NOT Increase!]
For an open room or vented attic of fixed volume $V$ exposed to constant atmospheric pressure $P_{\mathrm{atm}}$ (as in \texttt{2019 MIT E2 P1b}), heating the air from $T_1$ to $T_2$ causes warm air to escape so that $n(T) = P_{\mathrm{atm}}V/(RT)$. Consequently, the total internal energy of the air remaining inside the room is $U_{\mathrm{room}} = \frac{f}{2}nRT = \frac{f}{2}P_{\mathrm{atm}}V = \mathbf{\text{constant}}$!\idxerr{open-room heating constant internal energy trap}
\end{trapbox}
\end{twocol}

\begin{workedexample}[Two-Bulb Capillary Equilibrium \& Fractional Mass Redistribution]{Invitational Archetype \textbar\ Tier 2}{ex:06.2}
Two rigid glass bulbs of volumes $V_A = \SI{2.00}{\liter}$ and $V_B = \SI{3.00}{\liter}$ are connected by a negligible-volume stopcock. Initially, the stopcock is closed: Bulb~$A$ contains dry nitrogen ($\mathrm{N_2}$, $M = \SI{28.0134}{\gram\per\mole}$) at $P_{A,i} = \SI{250.0}{\kilo\pascal}$ and $T_{A,i} = \SI{300.0}{\kelvin}$, while Bulb~$B$ contains $\mathrm{N_2}$ at $P_{B,i} = \SI{100.0}{\kilo\pascal}$ and $T_{B,i} = \SI{400.0}{\kelvin}$. The stopcock is opened while Bulb~$A$ is maintained in a water bath at $T_{A,f} = \SI{300.0}{\kelvin}$ and Bulb~$B$ is heated to $T_{B,f} = \SI{500.0}{\kelvin}$. Find (a) the total moles $n_{\mathrm{tot}}$, (b) the final uniform pressure $P_f$, and (c) the net mass of $\mathrm{N_2}$ ($\Delta m_A$) that flows into or out of Bulb~$A$.
\par\smallskip\noindent
\textbf{Solution.}
\begin{enumerate}
  \item \textbf{(a) Initial and total moles (using $R = \SI{8.314463}{\liter\kilo\pascal\per\mole\per\kelvin}$):}
  \[
    n_{A,i} = \frac{P_{A,i}V_A}{RT_{A,i}} = \frac{250.0\times 2.00}{8.314463\times 300.0} = \SI{0.200454}{\mole}, \quad n_{B,i} = \frac{P_{B,i}V_B}{RT_{B,i}} = \frac{100.0\times 3.00}{8.314463\times 400.0} = \SI{0.090204}{\mole}
  \]
  \[
    n_{\mathrm{tot}} = n_{A,i} + n_{B,i} = \mathbf{\SI{0.290658}{\mole}} \quad (m_{\mathrm{tot}} = n_{\mathrm{tot}}M = \SI{8.1423}{\gram})
  \]
  \item \textbf{(b) Final pressure $P_f$ from \eref{eq:06:two-bulbs}:}
  \[
    P_f = \frac{n_{\mathrm{tot}}R}{\dfrac{V_A}{T_{A,f}} + \dfrac{V_B}{T_{B,f}}} = \frac{250.0(2.00)/300.0 + 100.0(3.00)/400.0}{2.00/300.0 + 3.00/500.0} = \frac{1.66667 + 0.75000}{0.0066667 + 0.0060000} = \frac{2.41667}{0.0126667} = \mathbf{\SI{190.789}{\kilo\pascal}}
  \]
  \item \textbf{(c) Final moles in Bulb $A$ and net mass transfer:}
  \[
    n_{A,f} = \frac{P_f V_A}{RT_{A,f}} = \frac{190.789\times 2.00}{8.314463\times 300.0} = \SI{0.152978}{\mole}, \qquad \Delta m_A = (n_{A,f} - n_{A,i})M = (-0.047476\,\si{\mole})(\SI{28.0134}{\gram\per\mole}) = \mathbf{\SI{-1.3300}{\gram}}
  \]
  Thus $\SI{1.330}{\gram}$ of $\mathrm{N_2}$ flows from Bulb~$A$ into Bulb~$B$.
\end{enumerate}
\verifiedby{Conservation check: $n_{B,f} = P_f V_B/(RT_{B,f}) = \SI{0.137680}{\mole}$ and $n_{A,f} + n_{B,f} = 0.152978 + 0.137680 = \SI{0.290658}{\mole} = n_{\mathrm{tot}}$ (\texttt{check\_sec06\_09.py::ex06\_2}).}
\end{workedexample}

% ============================================================================
\section{Dalton's Law of Partial Pressures, Wet Gases \& Graham's Law of Effusion}\label{sec:06.3}
% ============================================================================
\idxtopic{Dalton's law of partial pressures}\idxtopic{Amagat's law}\idxtopic{mole fraction}\idxtopic{Graham's law of effusion}\idxtopic{gas collection over water}

For a non-reacting mixture of $K$ ideal gases occupying volume $V$ at temperature $T$, each species $i$ behaves as if it alone occupied the entire volume \src{S13,S16}:
\begin{equation}\label{eq:06:dalton}
  P_{\mathrm{tot}} = \sum_{i=1}^K P_i = \left(\sum_{i=1}^K n_i\right)\frac{RT}{V} = n_{\mathrm{tot}}\frac{RT}{V}, \qquad P_i = x_i P_{\mathrm{tot}}, \qquad x_i = \frac{n_i}{n_{\mathrm{tot}}} = \frac{P_i}{P_{\mathrm{tot}}} = \frac{V_i}{V_{\mathrm{tot}}}
\end{equation}
\idxformula{Dalton's law: $P_{\mathrm{tot}} = \sum P_i$, $P_i = x_i P_{\mathrm{tot}}$}
where $V_i = x_i V_{\mathrm{tot}}$ is the \textbf{Amagat partial volume} (the volume species~$i$ would occupy alone at $P_{\mathrm{tot}}$ and $T$). The mixture's apparent molar mass $M_{\mathrm{mix}}$ and specific gas constant $R_{s,\mathrm{mix}}$ are:
\begin{equation}\label{eq:06:mmix}
  M_{\mathrm{mix}} = \frac{m_{\mathrm{tot}}}{n_{\mathrm{tot}}} = \sum_{i=1}^K x_i M_i = \left(\sum_{i=1}^K \frac{w_i}{M_i}\right)^{-1}, \qquad R_{s,\mathrm{mix}} = \frac{R}{M_{\mathrm{mix}}} = \sum_{i=1}^K w_i R_{s,i}
\end{equation}
\idxformula{Mixture molar mass: $M_{\mathrm{mix}} = \sum x_i M_i$}
where $w_i = m_i/m_{\mathrm{tot}} = x_i M_i / M_{\mathrm{mix}}$ is the mass fraction of species~$i$.

\begin{formulacard}[Wet-Gas Eudiometry \& Graham's Law of Pinhole Effusion \src{S09,S13}]
\begin{itemize}
  \item \textbf{Gas Collected Over Water (Eudiometer Tube Correction):} When a gas is bubbled through liquid water at temperature $T$ into an inverted graduated tube where the internal water level stands a height $h$ \emph{above} the outer reservoir level, the partial pressure of the dry gas is:
  \begin{equation}\label{eq:06:wet-gas}
    P_{\mathrm{tot}} = P_{\mathrm{atm}} - \rho_{\mathrm{H_2O}} g_n h = P_{\mathrm{dry}} + P_{\mathrm{sat,H_2O}}(T) \implies \boxed{P_{\mathrm{dry}} = P_{\mathrm{atm}} - \rho_{\mathrm{H_2O}} g_n h - P_{\mathrm{sat,H_2O}}(T)}
  \end{equation}
  \idxformula{Wet gas partial pressure: $P_{\mathrm{dry}} = P_{\mathrm{tot}} - P_{\mathrm{sat,H_2O}}(T)$}
  \item \textbf{Graham's Law of Effusion (1848):} The molecular flux escaping through a pinhole of area $A_h$ (diameter $\ll$ mean free path $\lambda$) into vacuum is $\dot{N} = \frac{1}{4}(N/V)v_{\mathrm{avg}}A_h = \frac{P A_h}{\sqrt{2\pi m_0 k_B T}}$. Comparing two gases at identical $T$ and $P$:
  \begin{equation}\label{eq:06:graham}
    \boxed{\frac{\dot{n}_1}{\dot{n}_2} = \frac{\dot{V}_1}{\dot{V}_2} = \frac{v_{\mathrm{avg},1}}{v_{\mathrm{avg},2}} = \sqrt{\frac{M_2}{M_1}}}, \qquad \frac{t_{\mathrm{eff},1}}{t_{\mathrm{eff},2}} = \sqrt{\frac{M_1}{M_2}}, \qquad \boxed{\frac{\dot{m}_1}{\dot{m}_2} = \frac{M_1\dot{n}_1}{M_2\dot{n}_2} = \sqrt{\frac{M_1}{M_2}}}
  \end{equation}
  \idxformula{Graham's law of effusion: $\dot{n}_1/\dot{n}_2 = \sqrt{M_2/M_1}$}
  \item \textbf{Multi-Stage Cascade Enrichment ($N_{\mathrm{stage}}$ porous-barrier stages):} If a binary mixture of initial molar ratio $(n_1/n_2)_0$ effuses sequentially through $N_{\mathrm{stage}}$ ideal vacuum stages, the enriched molar ratio is:
  \begin{equation}\label{eq:06:graham-cascade}
    \left(\frac{n_1}{n_2}\right)_{N_{\mathrm{stage}}} = \left(\frac{n_1}{n_2}\right)_0 \left(\frac{M_2}{M_1}\right)^{N_{\mathrm{stage}}/2} \implies N_{\mathrm{stage}} = \frac{2\ln\left[\dfrac{(n_1/n_2)_{N_{\mathrm{stage}}}}{(n_1/n_2)_0}\right]}{\ln(M_2/M_1)}
  \end{equation}
\end{itemize}
\end{formulacard}

\begin{workedexample}[Wet Oxygen Eudiometry \& Graham Effusion Mass-vs-Mole Trap]{Regional/State Standard \textbar\ Tier 2}{ex:06.3}
A $\SI{500.0}{\milli\liter}$ eudiometer tube collects oxygen gas ($\mathrm{O_2}$, $M = \SI{31.9988}{\gram\per\mole}$) over water at $T = \SI{25.0}{\celsius}$ ($\SI{298.15}{\kelvin}$, where $P_{\mathrm{sat,H_2O}} = \SI{3.1699}{\kilo\pascal}$ from NIST Table~3.1 \src{S09}). The water level inside the tube is $h = \SI{13.60}{\centi\meter}$ \emph{above} the outer dish ($\rho_{\mathrm{H_2O}} = \SI{997.0}{\kilogram\per\cubic\meter}$), and the barometer reads $P_{\mathrm{atm}} = \SI{101.325}{\kilo\pascal}$. (a) Find the dry mass of $\mathrm{O_2}$ collected. (b) If an equimolar mixture of $\mathrm{He}$ ($M_{\mathrm{He}} = \SI{4.0026}{\gram\per\mole}$) and $\mathrm{O_2}$ at $\SI{298.15}{\kelvin}$ effuses through a pinhole into vacuum, compute both the \textbf{molar} effusion ratio $\dot{n}_{\mathrm{He}}/\dot{n}_{\mathrm{O_2}}$ and the \textbf{mass} effusion ratio $\dot{m}_{\mathrm{He}}/\dot{m}_{\mathrm{O_2}}$.
\par\smallskip\noindent
\textbf{Solution.}
\begin{enumerate}
  \item \textbf{(a) Hydrostatic and vapor-pressure corrections from \eref{eq:06:wet-gas}:}
  \[
    \Delta P_{\mathrm{hydro}} = \rho_{\mathrm{H_2O}}g_n h = (997.0)(9.80665)(0.1360) = \SI{1329.7}{\pascal} = \SI{1.3297}{\kilo\pascal}
  \]
  \[
    P_{\mathrm{dry,O_2}} = 101.325 - 1.3297 - 3.1699 = \SI{96.8254}{\kilo\pascal}
  \]
  \[
    n_{\mathrm{O_2}} = \frac{P_{\mathrm{dry,O_2}}V}{RT} = \frac{(96.8254)(0.5000)}{(8.314463)(298.15)} = \SI{0.0195295}{\mole} \implies m_{\mathrm{O_2}} = n_{\mathrm{O_2}}M_{\mathrm{O_2}} = \mathbf{\SI{0.62492}{\gram}}
  \]
  \item \textbf{(b) Molar vs.\ mass effusion rate ratios from \eref{eq:06:graham}:}
  \[
    \frac{\dot{n}_{\mathrm{He}}}{\dot{n}_{\mathrm{O_2}}} = \sqrt{\frac{M_{\mathrm{O_2}}}{M_{\mathrm{He}}}} = \sqrt{\frac{31.9988}{4.0026}} = \mathbf{2.8275}, \qquad \frac{\dot{m}_{\mathrm{He}}}{\dot{m}_{\mathrm{O_2}}} = \sqrt{\frac{M_{\mathrm{He}}}{M_{\mathrm{O_2}}}} = \sqrt{\frac{4.0026}{31.9988}} = \mathbf{0.35367}
  \]
  \idxerr{Graham's law molar vs. mass effusion rate inversion}
  \emph{Trap check:} Helium effuses $2.83\times$ faster in \emph{moles per second}, but oxygen effuses $2.83\times$ faster in \emph{grams per second}!
\end{enumerate}
\verifiedby{Identity $(\dot{n}_{\mathrm{He}}/\dot{n}_{\mathrm{O_2}})(\dot{m}_{\mathrm{He}}/\dot{m}_{\mathrm{O_2}}) = 2.82746 \times 0.353674 = 1.00000$ (\texttt{check\_sec06\_09.py::ex06\_3}).}
\end{workedexample}

% ============================================================================
\section{Microscopic Kinetic Theory, Maxwell--Boltzmann Speeds \& Mean Free Path}\label{sec:06.4}
% ============================================================================
\idxtopic{root-mean-square speed}\idxtopic{most probable speed}\idxtopic{average speed}\idxtopic{speed of sound}\idxtopic{mean free path}\idxtopic{collision frequency}

\begin{derivation}[Algebra-Only Cubic-Box Collision Proof of Kinetic Theory Pressure \src{S13,S15}]
Consider $N$ point molecules of mass $m_0 = M/N_A$ moving without intermolecular forces inside a rigid cube of side $L$ and volume $V = L^3$:
\begin{enumerate}
  \item \textbf{Single-wall elastic impulse:} A molecule with velocity components $(v_x, v_y, v_z)$ strikes the right wall ($x = L$, area $A = L^2$) and rebounds elastically with $(-v_x, v_y, v_z)$, transferring momentum $\Delta p_x = +2m_0 v_x$ to the wall.
  \item \textbf{Round-trip collision interval:} The molecule travels distance $2L$ across the box and back before hitting the same wall again, requiring time $\Delta t = 2L/v_x$, so its time-averaged force on the wall is $f_x = \Delta p_x / \Delta t = m_0 v_x^2 / L$.
  \item \textbf{3D isotropy \& Pythagorean theorem:} Summing over all $N$ molecules and dividing by area $A = L^2$ gives $P = \frac{N m_0 \langle v_x^2\rangle}{V}$. Because thermal motion has no preferred axis, $\langle v_x^2\rangle = \langle v_y^2\rangle = \langle v_z^2\rangle = \frac{1}{3}\langle v_x^2 + v_y^2 + v_z^2\rangle = \frac{1}{3}\langle v^2\rangle = \frac{1}{3}v_{\mathrm{rms}}^2$, yielding:
  \begin{equation}\label{eq:06:kinetic-p}
    \boxed{P = \frac{1}{3}\frac{N}{V}m_0 v_{\mathrm{rms}}^2 = \frac{1}{3}\rho v_{\mathrm{rms}}^2 = \frac{2}{3}\left(\frac{N}{V}\right)K_{\mathrm{trans,avg}}}, \qquad \boxed{K_{\mathrm{trans,avg}} = \frac{1}{2}m_0 v_{\mathrm{rms}}^2 = \frac{3}{2}k_B T}
  \end{equation}
  \idxformula{Kinetic theory pressure: $P = \frac{1}{3}\rho v_{\mathrm{rms}}^2 = \frac{2}{3}(N/V)K_{\mathrm{trans,avg}}$}\idxformula{Average translational kinetic energy: $K_{\mathrm{trans,avg}} = \frac{3}{2}k_B T$}
\end{enumerate}
\end{derivation}

In thermal equilibrium at temperature $T$, James Clerk Maxwell (1860) and Ludwig Boltzmann (1868) showed that the fraction of molecules with speeds in $[v, v+\Delta v]$ follows the \textbf{Maxwell--Boltzmann speed distribution} \src{S13,S15,S23}:
\begin{equation}\label{eq:06:mb-dist}
  f(v) = 4\pi\left(\frac{m_0}{2\pi k_B T}\right)^{3/2} v^2 \exp\left(-\frac{m_0 v^2}{2k_B T}\right) = 4\pi\left(\frac{M}{2\pi R T}\right)^{3/2} v^2 \exp\left(-\frac{M v^2}{2 R T}\right)
\end{equation}
\idxformula{Maxwell--Boltzmann speed distribution: $f(v) = 4\pi(M/(2\pi RT))^{3/2}v^2 e^{-Mv^2/(2RT)}$}

\begin{formulacard}[Characteristic Molecular Speeds, Sound Speed \& Mean Free Path \src{S13,S15,S18}]
\par\smallskip\noindent
\footnotesize
\begin{tabularx}{\linewidth}{@{}l l l X@{}}
\toprule
\textbf{Characteristic Speed} & \textbf{Molar Form ($M$ in $\si{\kilogram\per\mole}$!)} & \textbf{Molecular Form ($m_0$ in $\si{\kilogram}$)} & \textbf{Physical Meaning \& Ratio to $v_{\mathrm{mp}}$} \\
\midrule
\textbf{Most Probable ($v_{\mathrm{mp}}$)} & $\displaystyle v_{\mathrm{mp}} = \sqrt{\frac{2RT}{M}}$ & $\displaystyle \sqrt{\frac{2k_B T}{m_0}}$ & Peak of $f(v)$ curve; $\mathbf{1.0000\,v_{\mathrm{mp}}}$ ($\sqrt{2} \approx 1.4142$) \\[6pt]
\textbf{Average ($v_{\mathrm{avg}} = \langle v\rangle$)} & $\displaystyle v_{\mathrm{avg}} = \sqrt{\frac{8RT}{\pi M}}$ & $\displaystyle \sqrt{\frac{8k_B T}{\pi m_0}}$ & Controls collision rate \& effusion; $\mathbf{1.1284\,v_{\mathrm{mp}}}$ ($\sqrt{8/\pi} \approx 1.5958$) \\[6pt]
\textbf{Root-Mean-Square ($v_{\mathrm{rms}}$)} & $\displaystyle v_{\mathrm{rms}} = \sqrt{\frac{3RT}{M}} = \sqrt{\frac{3P}{\rho}}$ & $\displaystyle \sqrt{\frac{3k_B T}{m_0}}$ & Controls $K_{\mathrm{trans}}$ \& pressure $P$; $\mathbf{1.2247\,v_{\mathrm{mp}}}$ ($\sqrt{3} \approx 1.7321$) \\[6pt]
\textbf{Adiabatic Sound ($c_s$)} & $\displaystyle c_s = \sqrt{\frac{\gamma RT}{M}} = \sqrt{\frac{\gamma P}{\rho}}$ & $\displaystyle \sqrt{\frac{\gamma k_B T}{m_0}} = \sqrt{\frac{\gamma}{3}}\,v_{\mathrm{rms}}$ & Acoustic wave speed ($\gamma = C_P/C_V$); for air ($\gamma=1.4$), $c_s = 0.6831\,v_{\mathrm{rms}}$ \\
\bottomrule
\end{tabularx}
\par\smallskip\noindent
\idxformula{Maxwell--Boltzmann speeds: $v_{\mathrm{mp}}=\sqrt{2RT/M} < v_{\mathrm{avg}}=\sqrt{8RT/(\pi M)} < v_{\mathrm{rms}}=\sqrt{3RT/M}$}\idxformula{Speed of sound in ideal gas: $c_s = \sqrt{\gamma RT/M}$}
\begin{itemize}
  \item \textbf{Clausius--Maxwell Hard-Sphere Mean Free Path ($\lambda$) and Collision Frequency ($Z_{\mathrm{coll}}$):} For molecules of effective collision diameter $d$ (typical $d \approx \SIrange{2.2e-10}{4.0e-10}{\meter}$; for $\mathrm{N_2}$, $d \approx \SI{3.70e-10}{\meter}$):
  \begin{equation}\label{eq:06:mfp}
    \boxed{\lambda = \frac{1}{\sqrt{2}\,\pi d^2 (N/V)} = \frac{k_B T}{\sqrt{2}\,\pi d^2 P}}, \qquad Z_{\mathrm{coll}} = \frac{v_{\mathrm{avg}}}{\lambda} = \sqrt{2}\,\pi d^2 \left(\frac{P}{k_B T}\right)\sqrt{\frac{8RT}{\pi M}} = \frac{4\pi d^2 P}{\sqrt{\pi m_0 k_B T}}
  \end{equation}
  \idxformula{Mean free path: $\lambda = k_B T / (\sqrt{2}\pi d^2 P)$}
  \item \textbf{Maxwell's Pressure-Independent Transport Law (1860/1866):} Because dynamic viscosity scales as $\mu \approx \frac{1}{3}\rho v_{\mathrm{avg}}\lambda$ and thermal conductivity scales as $k \approx \frac{1}{3}\rho c_v v_{\mathrm{avg}}\lambda$, the density $\rho \propto P$ and mean free path $\lambda \propto 1/P$ \textbf{cancel out exactly}! Thus dilute-gas $\mu$ and $k$ are \textbf{independent of pressure $P$} at fixed $T$ and scale as $\mu \propto \sqrt{T}$ and $k \propto \sqrt{T}$ \src{S12,S23}.
\end{itemize}
\end{formulacard}

\par\smallskip\noindent
\begin{minipage}{\linewidth}
\centering
\begin{tikzpicture}[font=\scriptsize]
\begin{axis}[
    width=0.92\linewidth,
    height=4.8cm,
    xlabel={Molecular speed $v$ ($\si{\meter\per\second}$)},
    ylabel={Probability density $f(v)$ ($10^{-3}\,\si{\second\per\meter}$)},
    xmin=0, xmax=2000,
    ymin=0, ymax=2.2,
    grid=both,
    grid style={gray!25},
    legend pos=north east,
    legend style={font=\scriptsize, fill=white, draw=navy},
    tick label style={font=\scriptsize},
    label style={font=\scriptsize\bfseries}
]
% N2 (M = 0.0280134 kg/mol) at T = 300 K: a = M/(2RT) = 0.0280134/(2*8.31446*300) = 5.6154e-6
% 4*pi*(a/pi)^(1.5) * 1e3 = 4*pi*(1.7874e-6)^(1.5)*1e3 = 3.0036e-5
\addplot[navy, very thick, domain=0:2000, samples=36]
  {3.0036e-5 * x^2 * exp(-5.6154e-6 * x^2)};
\addlegendentry{$\mathrm{N_2}$ at $T=\SI{300}{\kelvin}$ ($M=\SI{28.0}{\gram\per\mole}$)}

% N2 at T = 900 K: a = 5.6154e-6 / 3 = 1.8718e-6; prefactor = 3.0036e-5 / 3^1.5 = 5.7804e-6
\addplot[rust, very thick, dashed, domain=0:2000, samples=36]
  {5.7804e-6 * x^2 * exp(-1.8718e-6 * x^2)};
\addlegendentry{$\mathrm{N_2}$ at $T=\SI{900}{\kelvin}$ (broadened, peak $\div\sqrt{3}$)}

% Vertical markers for N2 at 300 K: v_mp = 422.1, v_avg = 476.3, v_rms = 516.9
\draw[navy, thick, dotted] (axis cs:422.1,0) -- (axis cs:422.1,1.97)
  node[above, fill=white, inner sep=1pt] {$v_{\mathrm{mp}}=\SI{422}{\meter\per\second}$};
\draw[teal, thick, dashdotted] (axis cs:476.3,0) -- (axis cs:476.3,1.89);
\node[teal, right, fill=white, inner sep=1pt] at (axis cs:485,1.58) {$v_{\mathrm{avg}}=\SI{476}{\meter\per\second}$};
\draw[black, thick] (axis cs:516.9,0) -- (axis cs:516.9,1.76);
\node[black, right, fill=white, inner sep=1pt] at (axis cs:525,1.30) {$v_{\mathrm{rms}}=\SI{517}{\meter\per\second}$};

% Vertical marker for N2 at 900 K: v_mp = 731.1, v_rms = 895.4
\draw[rust, thick, dotted] (axis cs:731.1,0) -- (axis cs:731.1,1.14)
  node[above, fill=white, inner sep=1pt] {$v_{\mathrm{mp}}(\SI{900}{\kelvin})=\SI{731}{\meter\per\second}$};
\draw[rust, thick] (axis cs:895.4,0) -- (axis cs:895.4,1.02)
  node[right, fill=white, inner sep=1pt] {$v_{\mathrm{rms}}(\SI{900}{\kelvin})=\SI{895}{\meter\per\second}$};
\end{axis}
\end{tikzpicture}
\captionof{figure}{Maxwell--Boltzmann molecular speed distribution $f(v)$ for $\mathrm{N_2}$ at $T = \SI{300}{\kelvin}$ and $T = \SI{900}{\kelvin}$ \src{S13,S15}. Tripling $T$ shifts all three characteristic speeds ($v_{\mathrm{mp}} < v_{\mathrm{avg}} < v_{\mathrm{rms}}$) upward by $\sqrt{3} \approx 1.732$ and lowers the peak height $f(v_{\mathrm{mp}}) \propto 1/\sqrt{T}$ by $\sqrt{3}$ so total probability area remains unity.}\label{fig:06:maxwell-boltzmann}
\idxfig{Maxwell--Boltzmann speed distribution at 300 K and 900 K}
\end{minipage}
\par\smallskip

\begin{workedexample}[Discrete Two-Speed Ensemble Trap vs.\ Mean Free Path in Dry Air]{2018 MIT E3 Q21 \textbar\ Tier 2}{ex:06.4}
(a) In a hypothetical ideal gas of $N$ particles (each of mass $m_0$), exactly half the particles move at speed $v_0$ and half move at speed $2v_0$. Express the kinetic temperature $T_{\mathrm{actual}}$ in terms of $m_0, v_0, k_B$, and compare it with the false temperature $T_{\mathrm{false}}$ obtained by substituting the mean speed $\langle v\rangle = 1.5v_0$ into $\frac{1}{2}m_0\langle v\rangle^2 = \frac{3}{2}k_B T$. (b) For dry nitrogen ($\mathrm{N_2}$, $M = \SI{0.0280134}{\kilogram\per\mole}$, collision diameter $d = \SI{3.70e-10}{\meter}$) at $T = \SI{300.0}{\kelvin}$ and $P = \SI{101.325}{\kilo\pascal}$, compute $v_{\mathrm{mp}}, v_{\mathrm{avg}}, v_{\mathrm{rms}}$, the mean free path $\lambda$, and the single-molecule collision frequency $Z_{\mathrm{coll}}$.
\par\smallskip\noindent
\textbf{Solution.}
\begin{enumerate}
  \item \textbf{(a) Discrete speed ensemble ($\langle v^2\rangle$ vs.\ $\langle v\rangle^2$ trap):}
  \[
    \langle v^2\rangle = \frac{1}{2}v_0^2 + \frac{1}{2}(2v_0)^2 = 2.50\,v_0^2 \implies T_{\mathrm{actual}} = \frac{m_0\langle v^2\rangle}{3k_B} = \mathbf{\frac{5m_0 v_0^2}{6k_B} = 0.8333\,\frac{m_0 v_0^2}{k_B}}
  \]
  Using $\langle v\rangle^2 = (1.5v_0)^2 = 2.25\,v_0^2$ yields $T_{\mathrm{false}} = \frac{2.25m_0 v_0^2}{3k_B} = \frac{3m_0 v_0^2}{4k_B} = 0.7500\,\frac{m_0 v_0^2}{k_B}$, underestimating $T_{\mathrm{actual}}$ by $10.0\%$ because $\mathrm{Var}(v) = \langle v^2\rangle - \langle v\rangle^2 = 0.25\,v_0^2 > 0$!\idxerr{substituting mean speed $\langle v\rangle$ instead of $v_{\mathrm{rms}}$ in temperature formula}
  \item \textbf{(b) Characteristic speeds and mean free path for $\mathrm{N_2}$ at $\SI{300.0}{\kelvin}, \SI{101.325}{\kilo\pascal}$:}
  \[
    v_{\mathrm{mp}} = \sqrt{\frac{2RT}{M}} = \mathbf{\SI{422.00}{\meter\per\second}}, \qquad v_{\mathrm{avg}} = \sqrt{\frac{8RT}{\pi M}} = \mathbf{\SI{476.17}{\meter\per\second}}, \qquad v_{\mathrm{rms}} = \sqrt{\frac{3(8.314463)(300.0)}{0.0280134}} = \mathbf{\SI{516.84}{\meter\per\second}}
  \]
  \[
    \lambda = \frac{k_B T}{\sqrt{2}\,\pi d^2 P} = \frac{(1.380649\times 10^{-23})(300.0)}{\sqrt{2}\,\pi (3.70\times 10^{-10})^2 (101325)} = \mathbf{\SI{6.721e-8}{\meter}} = \SI{67.21}{\nano\meter}, \qquad Z_{\mathrm{coll}} = \frac{v_{\mathrm{avg}}}{\lambda} = \mathbf{\SI{7.085e9}{\per\second}}
  \]
\end{enumerate}
\verifiedby{Number density $n_{\#} = P/(k_BT) = \SI{2.446e25}{\per\cubic\meter}$ gives $\lambda = 1/(\sqrt{2}\pi d^2 n_{\#}) = \SI{67.21}{\nano\meter}$ (\texttt{check\_sec06\_09.py::ex06\_4}).}
\end{workedexample}

% ============================================================================
\section{Equipartition of Energy, Mayer's Relation, Quantum Freeze-Out \& Gas Mixtures}\label{sec:06.5}
% ============================================================================
\idxtopic{equipartition theorem}\idxtopic{degrees of freedom}\idxtopic{Mayer's relation}\idxtopic{adiabatic index $\gamma$}\idxtopic{quantum freeze-out}\idxtopic{rotational temperature}\idxtopic{vibrational temperature}

By the classical \textbf{Equipartition Theorem} \src{S13,S15}, every independent quadratic term ($\frac{1}{2}mv_i^2$, $\frac{1}{2}I_j\omega_j^2$, or $\frac{1}{2}k_s q_k^2$) in a molecule's Hamiltonian contributes $\frac{1}{2}k_B T$ to the average molecular energy ($\frac{1}{2}RT$ per mole). Notice that a 1D harmonic vibrational mode has \textbf{two} quadratic terms (kinetic + potential), contributing $2\times\frac{1}{2}RT = RT$ per mole when fully excited! For a molecule of $N_{\mathrm{atom}}$ atoms, total modes are $3N_{\mathrm{atom}} = 3_{\mathrm{trans}} + f_{\mathrm{rot}} + f_{\mathrm{vib,modes}}$, where $f_{\mathrm{rot}} = 2$ (linear molecule, since moment of inertia about the internuclear axis is negligible) or $f_{\mathrm{rot}} = 3$ (nonlinear molecule).

\begin{formulacard}[Equipartition Relations, Mayer's Identity \& Quantum Mode Freeze-Out \src{S10,S13,S15}]
\begin{equation}\label{eq:06:equipartition}
  \boxed{U = \frac{f}{2}nRT = \frac{f}{2}PV}, \qquad \boxed{C_{V,m} = \frac{f}{2}R = \frac{R}{\gamma - 1}}, \qquad \boxed{C_{P,m} = C_{V,m} + R = \frac{f+2}{2}R = \frac{\gamma R}{\gamma - 1}}, \qquad \boxed{\gamma = \frac{C_P}{C_V} = 1 + \frac{2}{f}}
\end{equation}
\idxformula{Equipartition internal energy and heat capacities: $U = \frac{f}{2}nRT$, $C_{V,m} = \frac{f}{2}R$, $C_{P,m} = \frac{f+2}{2}R$, $\gamma = 1 + 2/f$}\idxformula{Mayer's relation: $C_{P,m} - C_{V,m} = R$}
\textbf{Algebra-Only Proof of Mayer's Relation ($C_{P,m} - C_{V,m} = R$):} By definition $H = U + PV$. For an ideal gas, $PV = nRT$ and $U = U(T)$ only (Joule's Second Law), so over any temperature change $\Delta T$: $\Delta H = \Delta U + \Delta(PV) \implies nC_{P,m}\Delta T = nC_{V,m}\Delta T + nR\Delta T \implies C_{P,m} - C_{V,m} = R$ (and $c_p - c_v = R_s = R/M$).
\par\smallskip\noindent
\textbf{Quantum Rotational \& Vibrational Freeze-Out (\texttt{2018 MIT E3 MC4} \& \texttt{2027 Cornell E1 Q35}):} A quantum mode of spacing $\Delta E = k_B\Theta$ is thermally active only when $T \gtrsim \Theta$. For $\mathrm{H_2}$, $\Theta_{\mathrm{rot}} = \frac{\hbar^2}{2I k_B} \approx \SI{85.4}{\kelvin}$ and $\Theta_{\mathrm{vib}} = \frac{\hbar\omega}{k_B} \approx \SI{6100}{\kelvin}$ (for $\mathrm{N_2}$, $\Theta_{\mathrm{rot}} \approx \SI{2.88}{\kelvin}$ and $\Theta_{\mathrm{vib}} \approx \SI{3374}{\kelvin}$):
\begin{itemize}
  \item \textbf{Cryogenic ($T \ll \Theta_{\mathrm{rot}}$, e.g.\ $\mathrm{H_2}$ below $\SI{50}{\kelvin}$):} Rotational and vibrational modes are frozen out $\implies f = 3$ ($C_{V,m} = \frac{3}{2}R, \gamma = \frac{5}{3}$).
  \item \textbf{Room Temperature ($\Theta_{\mathrm{rot}} \ll T \ll \Theta_{\mathrm{vib}}$, $\SI{300}{\kelvin}$):} Translation + rotation active, vibration frozen out $\implies f = 5$ ($C_{V,m} = \frac{5}{2}R, \gamma = \frac{7}{5} = 1.40$).
  \item \textbf{High Temperature ($T \gtrsim \Theta_{\mathrm{vib}}$, $\gtrsim \SI{3000}{\kelvin}$):} Translation + rotation + vibration (2 quadratic terms!) active $\implies f = 7$ ($C_{V,m} = \frac{7}{2}R, \gamma = \frac{9}{7} \approx 1.286$).
\end{itemize}
\end{formulacard}

\par\smallskip\noindent
\begin{minipage}{\linewidth}
\centering
\begin{tikzpicture}[font=\scriptsize]
\begin{axis}[
    width=0.92\linewidth,
    height=4.4cm,
    xlabel={$\log_{10}(T/\si{\kelvin})$ [Temperature $T$ from $\SI{10}{\kelvin}$ to $\SI{10000}{\kelvin}$]},
    ylabel={Dimensionless Molar Heat Capacity $C_{V,m}/R$},
    xmin=1.0, xmax=4.0,
    ymin=1.2, ymax=3.8,
    xtick={1.0, 1.5, 2.0, 2.477, 3.0, 3.5, 4.0},
    xticklabels={$1.0$ ($\SI{10}{\kelvin}$), $1.5$ ($\SI{32}{\kelvin}$), $2.0$ ($\SI{100}{\kelvin}$), $2.48$ ($\SI{300}{\kelvin}$), $3.0$ ($\SI{1000}{\kelvin}$), $3.5$ ($\SI{3162}{\kelvin}$), $4.0$ ($\SI{10000}{\kelvin}$)},
    ytick={1.5, 2.0, 2.5, 3.0, 3.5},
    yticklabels={$1.5$ ($f=3$), $2.0$, $2.5$ ($f=5$), $3.0$, $3.5$ ($f=7$)},
    grid=both,
    grid style={gray!25},
    tick label style={font=\scriptsize},
    label style={font=\scriptsize\bfseries}
]
% Smooth sigmoid staircase representing H2 rotational activation (~85 K -> log10=1.93) and vibrational activation (~2500 K -> log10=3.4)
\addplot[navy, very thick, domain=1.0:4.0, samples=36]
  {1.5 + 1.0/(1 + exp(-(x - 1.95)/0.14)) + 1.0/(1 + exp(-(x - 3.35)/0.18))};

\draw[rust, dashed, thick] (axis cs:1.0,1.5) -- (axis cs:4.0,1.5);
\draw[teal, dashed, thick] (axis cs:1.0,2.5) -- (axis cs:4.0,2.5);
\draw[amber, dashed, thick] (axis cs:1.0,3.5) -- (axis cs:4.0,3.5);
\draw[black, dotted, thick] (axis cs:2.477,1.2) -- (axis cs:2.477,3.8);

\node[fill=white, draw=navy, inner sep=1.5pt] at (axis cs:1.38,1.82) {\textbf{Trans.\ only} ($f=3$)};
\node[fill=white, draw=teal, inner sep=1.5pt] at (axis cs:2.48,2.82) {\textbf{Trans.\ + 2 Rot.} ($f=5$ at $\SI{300}{\kelvin}$)};
\node[fill=white, draw=amber, inner sep=1.5pt] at (axis cs:3.52,3.15) {\textbf{+ 1 Vib.\ (KE+PE)} ($f=7$)};
\node[rust, right, fill=white, inner sep=1pt] at (axis cs:3.18,1.85) {Dissociation $\gtrsim \SI{4500}{\kelvin}$};
\end{axis}
\end{tikzpicture}
\captionof{figure}{Molar heat capacity $C_{V,m}/R$ of diatomic hydrogen ($\mathrm{H_2}$) vs.\ $\log_{10}(T/\si{\kelvin})$ showing quantum freeze-out of vibrational modes below $\sim\SI{1000}{\kelvin}$ and rotational modes below $\sim\SI{85}{\kelvin}$ \src{S10,S13,S15}.}\label{fig:06:equipartition}
\idxfig{Equipartition staircase and quantum rotational/vibrational freeze-out of H2}
\end{minipage}
\par\smallskip

\begin{table}[!htbp]
\centering
\captionof{table}{Equipartition Ideal Limits vs.\ NIST Measured Gas Properties at $T = \SI{298.15}{\kelvin}$ ($\SI{25}{\celsius}$) and $P = \SI{1}{\text{atm}}$ ($\SI{101.325}{\kilo\pascal}$) \src{S10,S12,S13}.}\label{tab:06:gas-properties}
\footnotesize\setlength{\tabcolsep}{3.2pt}\renewcommand{\arraystretch}{1.10}
\begin{tabularx}{\linewidth}{@{}l l r r c r r r r r r r r r@{}}
\toprule
\textbf{Gas} & \textbf{Formula} & $M$ & $R_s$ & $f_{\mathrm{id}}$ & $C_{V,\mathrm{id}}/R$ & $\gamma_{\mathrm{id}}$ & \textbf{NIST} & \textbf{NIST} $C_{V,m}$ & \textbf{NIST} $C_{P,m}$ & \textbf{NIST} $c_v$ & \textbf{NIST} $c_p$ & \textbf{NIST} $\gamma$ & $c_s$ \\
 & & ($\si{\gram\per\mole}$) & ($\si{\joule\per\kilogram\per\kelvin}$) & & & & $C_{V,m}/R$ & ($\si{\joule\per\mole\per\kelvin}$) & ($\si{\joule\per\mole\per\kelvin}$) & ($\si{\joule\per\kilogram\per\kelvin}$) & ($\si{\joule\per\kilogram\per\kelvin}$) & ($\SI{25}{\celsius}$) & ($\si{\meter\per\second}$) \\
\midrule
Helium & $\mathrm{He}$ & $4.0026$ & $2077.26$ & $3$ & $1.50$ & $5/3$ & $1.500$ & $12.47$ & $20.79$ & $3116.2$ & $5193.1$ & $1.6665$ & $1016.4$ \\
Neon & $\mathrm{Ne}$ & $20.1797$ & $412.02$ & $3$ & $1.50$ & $5/3$ & $1.500$ & $12.47$ & $20.79$ & $618.1$ & $1030.3$ & $1.6670$ & $452.7$ \\
Argon & $\mathrm{Ar}$ & $39.9480$ & $208.13$ & $3$ & $1.50$ & $5/3$ & $1.501$ & $12.48$ & $20.84$ & $312.4$ & $521.6$ & $1.6696$ & $321.7$ \\
Hydrogen & $\mathrm{H_2}$ & $2.0159$ & $4124.48$ & $5$ & $2.50$ & $7/5$ & $2.468$ & $20.52$ & $28.84$ & $10179.7$ & $14306.4$ & $1.4054$ & $1315.4$ \\
Nitrogen & $\mathrm{N_2}$ & $28.0134$ & $296.80$ & $5$ & $2.50$ & $7/5$ & $2.504$ & $20.82$ & $29.17$ & $743.1$ & $1041.4$ & $1.4013$ & $352.1$ \\
Oxygen & $\mathrm{O_2}$ & $31.9988$ & $259.84$ & $5$ & $2.50$ & $7/5$ & $2.534$ & $21.07$ & $29.43$ & $658.5$ & $919.6$ & $1.3966$ & $328.7$ \\
Dry Air & $\mathrm{Air}$ & $28.9655$ & $287.05$ & $5$ & $2.50$ & $7/5$ & $2.503$ & $20.82$ & $29.14$ & $718.6$ & $1006.2$ & $1.4002$ & $346.1$ \\
Carbon monoxide & $\mathrm{CO}$ & $28.0101$ & $296.84$ & $5$ & $2.50$ & $7/5$ & $2.505$ & $20.83$ & $29.19$ & $743.7$ & $1042.2$ & $1.4013$ & $352.0$ \\
Carbon dioxide$^*$ & $\mathrm{CO_2}$ & $44.0095$ & $188.92$ & $5$ & $2.50$ & $7/5$ & $\mathbf{3.480}$ & $28.94$ & $37.45$ & $657.5$ & $850.9$ & $1.2941$ & $268.6$ \\
Water vapor ($\SI{100}{\celsius}$) & $\mathrm{H_2O}(g)$ & $18.0153$ & $461.52$ & $6$ & $3.00$ & $4/3$ & $3.371$ & $28.02$ & $37.47$ & $1555.6$ & $2079.8$ & $1.3370$ & $472.2$ \\
Methane & $\mathrm{CH_4}$ & $16.0425$ & $518.28$ & $6$ & $3.00$ & $4/3$ & $3.297$ & $27.41$ & $35.80$ & $1708.6$ & $2231.7$ & $1.3062$ & $448.5$ \\
Ammonia & $\mathrm{NH_3}$ & $17.0305$ & $488.21$ & $6$ & $3.00$ & $4/3$ & $3.366$ & $27.98$ & $36.84$ & $1643.2$ & $2163.2$ & $1.3165$ & $433.1$ \\
Sulfur dioxide & $\mathrm{SO_2}$ & $64.0638$ & $129.78$ & $6$ & $3.00$ & $4/3$ & $3.905$ & $32.47$ & $41.60$ & $506.8$ & $649.4$ & $1.2813$ & $218.9$ \\
Sulfur hexafluoride & $\mathrm{SF_6}$ & $146.0554$ & $56.93$ & $6$ & $3.00$ & $4/3$ & $\mathbf{10.699}$ & $88.96$ & $97.71$ & $609.1$ & $669.0$ & $1.0984$ & $135.0$ \\
\bottomrule
\end{tabularx}
\par\vspace{1pt}\raggedright{\scriptsize $^*$Linear triatomic $\mathrm{CO_2}$ has two degenerate low-frequency bending modes ($\Theta_{\mathrm{bend}} = \SI{960}{\kelvin}$) partially excited at $\SI{298}{\kelvin}$, elevating $C_{V,m}/R$ from $2.50$ to $3.48$ (\texttt{2019 MIT E2 MC13}) \src{S13,S26}.}
\end{table}

\begin{workedexample}[Adiabatic Index \& Sound Speed of a Multi-Component Gas Mixture]{State/National Standard \textbar\ Tier 2}{ex:06.5}
A $\mathrm{CO_2}$ gas-laser cavity at $T = \SI{300.0}{\kelvin}$ contains a mixture of $n_1 = \SI{0.700}{\mole}$ Helium ($\mathrm{He}$, monatomic $f_1=3$, $M_1 = \SI{4.0026}{\gram\per\mole}$), $n_2 = \SI{0.200}{\mole}$ Nitrogen ($\mathrm{N_2}$, rigid diatomic $f_2=5$, $M_2 = \SI{28.0134}{\gram\per\mole}$), and $n_3 = \SI{0.100}{\mole}$ Carbon Dioxide ($\mathrm{CO_2}$, using NIST $C_{V,m,3} = 3.480\,R$, $M_3 = \SI{44.0095}{\gram\per\mole}$). Find (a) the mixture molar heat capacities $C_{V,m,\mathrm{mix}}$ and $C_{P,m,\mathrm{mix}}$, (b) the true mixture adiabatic index $\gamma_{\mathrm{mix}}$ (and compare with the wrong mole-weighted average of $\gamma_i$), and (c) the speed of sound $c_{s,\mathrm{mix}}$ at $\SI{300.0}{\kelvin}$.
\par\smallskip\noindent
\textbf{Solution.}
\begin{enumerate}
  \item \textbf{(a) Mole-weighted extensive heat capacity average (NEVER average $\gamma_i$ directly!):\idxerr{averaging $\gamma_i$ directly in a gas mixture instead of $C_{P,\mathrm{mix}}/C_{V,\mathrm{mix}}$}}
  With $n_{\mathrm{tot}} = 1.000\,\si{\mole}$, $x_1 = 0.700, x_2 = 0.200, x_3 = 0.100$:
  \[
    \frac{C_{V,m,\mathrm{mix}}}{R} = 0.700(1.500) + 0.200(2.500) + 0.100(3.480) = 1.050 + 0.500 + 0.348 = \mathbf{1.8980} \implies C_{V,m,\mathrm{mix}} = \mathbf{\SI{15.781}{\joule\per\mole\per\kelvin}}
  \]
  \[
    C_{P,m,\mathrm{mix}} = C_{V,m,\mathrm{mix}} + R = 2.8980\,R = \mathbf{\SI{24.095}{\joule\per\mole\per\kelvin}}
  \]
  \item \textbf{(b) Mixture adiabatic index $\gamma_{\mathrm{mix}}$:}
  \[
    \gamma_{\mathrm{mix}} = \frac{C_{P,m,\mathrm{mix}}}{C_{V,m,\mathrm{mix}}} = \frac{2.8980}{1.8980} = \mathbf{1.5269}
  \]
  (Averaging $\gamma_i$ directly gives $0.70(1.6667) + 0.20(1.4000) + 0.10(1.2874) = 1.5754$, an error of $+3.18\%$!)
  \item \textbf{(c) Mixture molar mass and speed of sound:}
  \[
    M_{\mathrm{mix}} = 0.700(4.0026) + 0.200(28.0134) + 0.100(44.0095) = \SI{12.8055}{\gram\per\mole} = \SI{0.0128055}{\kilogram\per\mole}
  \]
  \[
    c_{s,\mathrm{mix}} = \sqrt{\frac{\gamma_{\mathrm{mix}} R T}{M_{\mathrm{mix}}}} = \sqrt{\frac{(1.52687)(8.314463)(300.0)}{0.0128055}} = \mathbf{\SI{545.36}{\meter\per\second}}
  \]
\end{enumerate}
\verifiedby{Direct check $C_{P,m,\mathrm{mix}} - C_{V,m,\mathrm{mix}} = 24.0953 - 15.7809 = \SI{8.3145}{\joule\per\mole\per\kelvin} = R$ (\texttt{check\_sec06\_09.py::ex06\_5}).}
\end{workedexample}

% ============================================================================
\section{Real-Gas Deviations: Compressibility Factor $Z$ \& Van der Waals Equation}\label{sec:06.6}
% ============================================================================
\idxtopic{compressibility factor $Z$}\idxtopic{van der Waals equation}\idxtopic{Boyle temperature}\idxtopic{critical point}\idxtopic{law of corresponding states}

Real gases deviate from $PV = nRT$ at high pressures (small intermolecular spacing) and low temperatures (near condensation) \src{S10,S13,S16}. The deviation is quantified by the dimensionless \textbf{compressibility factor $Z$} and the \textbf{virial expansion}:
\begin{equation}\label{eq:06:z-virial}
  Z \equiv \frac{PV}{nRT} = \frac{P V_m}{RT} = 1 + \frac{B(T)}{V_m} + \frac{C(T)}{V_m^2} + \cdots \approx 1 + \frac{1}{RT}\left(b - \frac{a}{RT}\right)P
\end{equation}
\idxformula{Compressibility factor: $Z = PV/(nRT) = PV_m/(RT)$}
\begin{itemize}
  \item \textbf{Moderate Pressure ($Z < 1$, Attractive Regime):} Long-range intermolecular attractions ($a > 0$: London dispersion, dipole--dipole, hydrogen bonding) reduce the impact momentum of molecules hitting container walls, lowering $P < P_{\mathrm{ideal}}$ and making the gas \emph{more compressible} than an ideal gas ($B(T) < 0$ for $T < T_B$).
  \item \textbf{Very High Pressure or High Temperature ($Z > 1$, Repulsive Excluded-Volume Regime):} Short-range hard-core repulsion ($b \approx 4N_A V_{\mathrm{molecule}} > 0$) reduces the free volume available for translation to $V - nb$, raising $P > P_{\mathrm{ideal}}$. For light gases with tiny $a$ ($\mathrm{He}, \mathrm{H_2}, \mathrm{Ne}$), $T_B < \SI{130}{\kelvin}$, so $Z > 1$ at all pressures at room temperature ($\SI{300}{\kelvin}$)!
\end{itemize}

\begin{formulacard}[Van der Waals Equation of State (1873), Critical Parameters \& Internal Energy \src{S10,S13,S16}]
\begin{equation}\label{eq:06:vdw}
  \boxed{\left(P + \frac{a n^2}{V^2}\right)(V - nb) = nRT} \iff \boxed{P = \frac{RT}{V_m - b} - \frac{a}{V_m^2}}, \qquad Z_{\mathrm{vdW}} = \frac{V_m}{V_m - b} - \frac{a}{RTV_m}
\end{equation}
\idxformula{Van der Waals equation: $(P + an^2/V^2)(V - nb) = nRT$}
\begin{itemize}
  \item \textbf{Critical Point Inflection Identity:} At the critical point $(T_c, P_c, V_{c,m})$, the isotherm has a horizontal inflection point, yielding the exact algebraic relations between $(a, b)$ and $(T_c, P_c, V_{c,m})$:
  \begin{equation}\label{eq:06:vdw-crit}
    V_{c,m} = 3b, \quad P_c = \frac{a}{27b^2}, \quad T_c = \frac{8a}{27Rb} \iff \boxed{a = \frac{27 R^2 T_c^2}{64 P_c}, \quad b = \frac{R T_c}{8 P_c}, \quad Z_{c,\mathrm{vdW}} = \frac{P_c V_{c,m}}{R T_c} = \frac{3}{8} = 0.375}
  \end{equation}
  \idxformula{Van der Waals critical constants: $a = 27R^2T_c^2/(64P_c)$, $b = RT_c/(8P_c)$, $Z_c = 3/8$}
  \item \textbf{Boyle Temperature ($T_B$, where $B(T_B) = 0$) \& Max Joule--Thomson Inversion Temperature ($T_{\mathrm{inv,max}}$):}
  \begin{equation}\label{eq:06:boyle-inv}
    T_B = \frac{a}{Rb} = \frac{27}{8}T_c = 3.375\,T_c, \qquad T_{\mathrm{inv,max}} = \frac{2a}{Rb} = 2T_B = \frac{27}{4}T_c = 6.75\,T_c
  \end{equation}
  \item \textbf{Law of Corresponding States:} Defining reduced variables $P_r = P/P_c$, $T_r = T/T_c$, and $V_r = V_m/V_{c,m}$ eliminates $(a,b)$ completely:
  \begin{equation}\label{eq:06:reduced-vdw}
    \left(P_r + \frac{3}{V_r^2}\right)(3V_r - 1) = 8T_r
  \end{equation}
  \item \textbf{Van der Waals Internal Energy \& Free-Expansion Cooling:} Because work must be done against attractive forces $a n^2/V^2$ when separating molecules, the internal energy of a Van der Waals gas depends on volume $V$ as well as $T$:
  \begin{equation}\label{eq:06:vdw-u}
    \boxed{U_{\mathrm{vdW}}(T,V) = nC_{V,m}T - \frac{a n^2}{V}} \implies \Delta T_{\text{free exp }(\Delta U = 0)} = -\frac{an}{C_{V,m}}\left(\frac{1}{V_1} - \frac{1}{V_2}\right) < 0 \quad (\text{always cools!})
  \end{equation}
  \idxformula{Van der Waals internal energy: $U_{\mathrm{vdW}} = nC_{V,m}T - an^2/V$}
\end{itemize}
\end{formulacard}

\par\smallskip\noindent
\begin{minipage}{\linewidth}
\centering
\begin{tikzpicture}[font=\scriptsize]
% Left Panel: Z vs P (bar) at 300 K
\begin{axis}[
    name=leftplot,
    width=0.47\linewidth,
    height=4.3cm,
    xlabel={Pressure $P$ ($\si{\bar}$)},
    ylabel={Compressibility $Z = PV/(nRT)$},
    xmin=0, xmax=400,
    ymin=0.3, ymax=1.3,
    grid=both,
    grid style={gray!25},
    legend pos=north west,
    legend style={font=\scriptsize, fill=white, draw=navy, inner sep=1pt},
    tick label style={font=\scriptsize},
    label style={font=\scriptsize\bfseries},
    title={\textbf{(a) Real-gas $Z(P)$ at $\SI{320}{\kelvin}$ (Peng--Robinson)}}
]
\draw[black, dashed, thick] (axis cs:0,1.0) -- (axis cs:400,1.0);
\node[font=\scriptsize, fill=white, inner sep=1pt] at (axis cs:218,0.94) {Ideal ($Z=1$)};
% H2 (always Z > 1 at 320 K)
\addplot[teal, very thick, smooth] coordinates {(0.01,1.000) (25,1.007) (50,1.016) (75,1.025) (100,1.034) (150,1.055) (200,1.078) (250,1.103) (300,1.128) (350,1.155) (400,1.183)};
\addlegendentry{$\mathrm{H_2}$ ($T > T_B$)}
% N2 (shallow dip then > 1)
\addplot[navy, very thick, densely dotted, smooth] coordinates {(0.01,1.000) (25,0.995) (50,0.994) (75,0.996) (100,1.000) (150,1.017) (200,1.041) (250,1.071) (300,1.105) (350,1.142) (400,1.182)};
\addlegendentry{$\mathrm{N_2}$}
% CH4 (deeper dip)
\addplot[amber, very thick, dashed, smooth] coordinates {(0.01,1.000) (25,0.959) (50,0.923) (75,0.893) (100,0.871) (150,0.849) (200,0.853) (250,0.876) (300,0.912) (350,0.955) (400,1.003)};
\addlegendentry{$\mathrm{CH_4}$}
% CO2 just above T_c = 304 K: deep dip near 120 bar (at 300 K it would condense near 67 bar instead)
\addplot[rust, very thick, smooth] coordinates {(0.01,1.000) (10,0.955) (20,0.909) (30,0.860) (40,0.809) (50,0.754) (60,0.695) (70,0.630) (80,0.556) (90,0.471) (100,0.391) (110,0.353) (120,0.344) (130,0.346) (140,0.353) (150,0.362) (175,0.390) (200,0.422) (225,0.455) (250,0.488) (275,0.521) (300,0.554) (325,0.587) (350,0.620) (375,0.653) (400,0.686)};
\addlegendentry{$\mathrm{CO_2}$ ($T_c = \SI{304}{\kelvin}$)}
\node[rust, fill=white, inner sep=1pt] at (axis cs:270,0.355) {$Z<1$: attraction ($a$)};
\node[teal, fill=white, inner sep=1pt] at (axis cs:300,1.24) {$Z>1$: molecular size ($b$)};
\end{axis}

% Right Panel: Reduced vdW Isotherms Pr vs Vr
\begin{axis}[
    at={(leftplot.outer east)}, anchor=outer west, xshift=0.25cm,
    width=0.47\linewidth,
    height=4.3cm,
    xlabel={Reduced molar volume $V_r = V_m/V_{c,m}$},
    ylabel={Reduced pressure $P_r = P/P_c$},
    xmin=0.45, xmax=3.0,
    ymin=0.0, ymax=2.0,
    grid=both,
    grid style={gray!25},
    legend pos=north east,
    legend style={font=\scriptsize, fill=white, draw=navy, inner sep=1pt},
    tick label style={font=\scriptsize},
    label style={font=\scriptsize\bfseries},
    title={\textbf{(b) Van der Waals Reduced Isotherms $P_r(V_r)$}}
]
% Tr = 1.15 (supercritical)
\addplot[teal, very thick, domain=0.55:3.0, samples=36] {8*1.15/(3*x - 1) - 3/(x^2)};
\addlegendentry{$T_r = 1.15 > 1$}
% Tr = 1.00 (critical isotherm through (1,1))
\addplot[navy, very thick, domain=0.52:3.0, samples=36] {8*1.00/(3*x - 1) - 3/(x^2)};
\addlegendentry{$T_r = 1.00$ (Critical)}
% Tr = 0.88 (subcritical spinodal loop + Maxwell equal-area tie-line at Pr ~ 0.58)
\addplot[rust, very thick, dashed, domain=0.50:3.0, samples=36] {8*0.88/(3*x - 1) - 3/(x^2)};
\addlegendentry{$T_r = 0.88 < 1$}
\draw[rust, ultra thick] (axis cs:0.575,0.579) -- (axis cs:2.04,0.579)
  node[pos=0.7, below, fill=white, inner sep=1pt] {Maxwell tie-line};
\fill[navy] (axis cs:1.0,1.0) circle (2pt) node[above right, fill=white, inner sep=1pt] {CP $(1,1)$};
\end{axis}
\end{tikzpicture}
\captionof{figure}{(a) Real-gas compressibility factor $Z = PV/(nRT)$ at $\SI{320}{\kelvin}$ (Peng--Robinson equation of state) showing attractive depression ($Z<1$) at moderate pressure and excluded-volume rise ($Z>1$) at high pressure. $\SI{320}{\kelvin}$ is chosen because $\mathrm{CO_2}$ is then supercritical; at $\SI{300}{\kelvin} < T_c$ it would liquefy near $\SI{67}{bar}$. (b) Reduced Van der Waals isotherms \eref{eq:06:reduced-vdw} showing the critical inflection point $(V_r, P_r)=(1,1)$ and Maxwell's equal-area vapor--liquid tie-line for $T_r < 1$ \src{S13,S16}.}\label{fig:06:vdw-z}
\idxfig{Real-gas compressibility factor Z and reduced van der Waals isotherms}
\end{minipage}
\par\smallskip

\begin{table}[!htbp]
\centering
\captionof{table}{Triple Points, Normal Boiling Points, Critical Constants \& Van der Waals Parameters ($a, b$) \src{S09,S10}.}\label{tab:06:vdw-constants}
\footnotesize\setlength{\tabcolsep}{3.5pt}\renewcommand{\arraystretch}{1.08}
\begin{tabularx}{\linewidth}{@{}l r r r r r r r r r r r r@{}}
\toprule
\textbf{Gas} & $T_{\mathrm{tp}}$ & $P_{\mathrm{tp}}$ & $T_b$ ($\SI{1}{\text{atm}}$) & $T_c$ & $P_c$ & $\rho_c$ & $Z_{c,\mathrm{exp}}$ & $a_{\mathrm{crit}}$ & $b_{\mathrm{crit}}$ & $a_{\mathrm{CRC}}$ & $b_{\mathrm{CRC}}$ & $T_{B,\mathrm{vdW}}$ \\
 & ($\si{\kelvin}$) & ($\si{\kilo\pascal}$) & ($\si{\kelvin}$) & ($\si{\kelvin}$) & ($\si{\kilo\pascal}$) & ($\si{\kilogram\per\cubic\meter}$) & & ($\si{\pascal\meter^6\per\mole^2}$) & ($10^{-5}\si{\cubic\meter\per\mole}$) & ($\si{\liter^2\bar\per\mole^2}$) & ($\si{\liter\per\mole}$) & ($\si{\kelvin}$) \\
\midrule
$\mathrm{He}$ & $2.18$ & $5.04$ & $4.22$ & $5.20$ & $228.3$ & $69.58$ & $0.304$ & $0.0034$ & $2.365$ & $0.0346$ & $0.0238$ & $17.5$ \\
$\mathrm{Ne}$ & $24.56$ & $43.36$ & $27.10$ & $44.40$ & $2661.6$ & $486.00$ & $0.299$ & $0.0216$ & $1.734$ & $0.2135$ & $0.0171$ & $149.9$ \\
$\mathrm{Ar}$ & $83.81$ & $68.89$ & $87.30$ & $150.69$ & $4863.0$ & $535.60$ & $0.290$ & $0.1362$ & $3.220$ & $1.3550$ & $0.0320$ & $508.6$ \\
$\mathrm{H_2}$ & $13.96$ & $7.36$ & $20.37$ & $33.15$ & $1296.4$ & $31.26$ & $0.303$ & $0.0247$ & $2.657$ & $0.2476$ & $0.0266$ & $111.9$ \\
$\mathrm{N_2}$ & $63.15$ & $12.52$ & $77.36$ & $126.19$ & $3395.8$ & $313.30$ & $0.289$ & $0.1368$ & $3.862$ & $1.3700$ & $0.0387$ & $425.9$ \\
$\mathrm{O_2}$ & $54.36$ & $0.15$ & $90.19$ & $154.58$ & $5043.0$ & $436.10$ & $0.288$ & $0.1382$ & $3.186$ & $1.3820$ & $0.0319$ & $521.7$ \\
$\mathrm{Air}$ & $59.75$ & $5.27$ & $78.90$ & $132.53$ & $3786.0$ & $342.68$ & $0.289$ & $0.1353$ & $3.638$ & $1.3580$ & $0.0364$ & $447.3$ \\
$\mathrm{CO}$ & $68.16$ & $15.54$ & $81.64$ & $132.86$ & $3494.0$ & $303.90$ & $0.292$ & $0.1473$ & $3.952$ & $1.4720$ & $0.0395$ & $448.4$ \\
$\mathrm{CO_2}$ & $216.59$ & $517.96$ & $194.69^*$ & $304.13$ & $7377.3$ & $467.60$ & $0.275$ & $0.3656$ & $4.285$ & $3.6580$ & $0.0429$ & $1026.4$ \\
$\mathrm{H_2O}$ & $273.16$ & $0.6117$ & $373.12$ & $647.10$ & $22064.0$ & $322.00$ & $0.229$ & $0.5535$ & $3.048$ & $5.5370$ & $0.0305$ & $2184.0$ \\
$\mathrm{CH_4}$ & $90.69$ & $11.70$ & $111.67$ & $190.56$ & $4599.2$ & $162.66$ & $0.286$ & $0.2303$ & $4.306$ & $2.3000$ & $0.0430$ & $643.1$ \\
$\mathrm{NH_3}$ & $195.49$ & $6.05$ & $239.83$ & $405.56$ & $11363.4$ & $233.25$ & $0.246$ & $0.4221$ & $3.709$ & $4.2250$ & $0.0371$ & $1368.8$ \\
$\mathrm{SO_2}$ & $197.70$ & $1.67$ & $263.14$ & $430.64$ & $7886.6$ & $517.50$ & $0.273$ & $0.6858$ & $5.675$ & $6.8650$ & $0.0568$ & $1453.4$ \\
$\mathrm{SF_6}$ & $223.56$ & $231.42$ & $204.90^*$ & $318.72$ & $3755.0$ & $742.30$ & $0.279$ & $0.7890$ & $8.822$ & $7.8570$ & $0.0879$ & $1075.7$ \\
$\mathrm{R134a}$ & $169.85$ & $0.39$ & $247.08$ & $374.21$ & $4059.3$ & $511.90$ & $0.260$ & $1.0061$ & $9.581$ & $10.090$ & $0.0957$ & $1263.0$ \\
\bottomrule
\end{tabularx}
\par\vspace{1pt}\raggedright{\scriptsize $^*$At $\SI{1}{\text{atm}}$, $\mathrm{CO_2}$ and $\mathrm{SF_6}$ sublime directly from solid to gas ($P_{\mathrm{tp}} > \SI{101.325}{\kilo\pascal}$). Note unit conversion: $a[\si{\liter^2\bar\per\mole^2}] = 10\,a[\si{\pascal\meter^6\per\mole^2}]$; $1\,\si{\liter^2\text{atm}\per\mole^2} = 1.01325\,\si{\liter^2\bar\per\mole^2}$.}
\end{table}

\begin{workedexample}[Van der Waals Pressure, Compressibility $Z$ \& Free-Expansion Cooling of $\mathrm{CO_2}$]{National/Invitational Standard \textbar\ Tier 3}{ex:06.6}
Consider $n = \SI{10.00}{\mole}$ of carbon dioxide ($\mathrm{CO_2}$, $a = \SI{0.3658}{\pascal\meter^6\per\mole^2} = \SI{3.658}{\liter^2\bar\per\mole^2}$, $b = \SI{4.286e-5}{\cubic\meter\per\mole} = \SI{0.04286}{\liter\per\mole}$, $C_{V,m} = \SI{28.94}{\joule\per\mole\per\kelvin}$) compressed into a rigid vessel of volume $V_1 = \SI{5.000}{\liter}$ ($\SI{5.000e-3}{\cubic\meter}$) at $T_1 = \SI{350.0}{\kelvin}$. (a) Compare the ideal-gas pressure $P_{\mathrm{ideal}}$ with the Van der Waals pressure $P_{\mathrm{vdW}}$ and find $Z_{\mathrm{vdW}}$. (b) If a partition ruptures and the gas undergoes an adiabatic free expansion into an evacuated chamber to final total volume $V_2 = \SI{50.00}{\liter}$, calculate the temperature drop $\Delta T$ from \eref{eq:06:vdw-u}.
\par\smallskip\noindent
\textbf{Solution.}
\begin{enumerate}
  \item \textbf{(a) Ideal vs.\ Van der Waals pressure at $V_m = V_1/n = \SI{0.5000}{\liter\per\mole} = \SI{5.000e-4}{\cubic\meter\per\mole}$:}
  \[
    P_{\mathrm{ideal}} = \frac{RT_1}{V_m} = \frac{(8.314463)(350.0)}{5.000\times 10^{-4}} = \SI{5.8201e6}{\pascal} = \mathbf{\SI{58.201}{\bar}}
  \]
  \[
    P_{\mathrm{vdW}} = \frac{RT_1}{V_m - b} - \frac{a}{V_m^2} = \frac{(8.314463)(350.0)}{(5.000 - 0.4286)\times 10^{-4}} - \frac{0.3658}{(5.000\times 10^{-4})^2} = 6.3658\times 10^6 - 1.4632\times 10^6 = \SI{4.9026e6}{\pascal} = \mathbf{\SI{49.026}{\bar}}
  \]
  \[
    Z_{\mathrm{vdW}} = \frac{P_{\mathrm{vdW}}}{P_{\mathrm{ideal}}} = \frac{49.026}{58.201} = \mathbf{0.8424} < 1 \quad (\text{attractions reduce } P \text{ by } \SI{9.18}{\bar} \text{ or } 15.8\%!)
  \]
  \item \textbf{(b) Adiabatic free-expansion cooling ($Q = 0, W_{\mathrm{by}} = 0 \implies \Delta U_{\mathrm{vdW}} = 0$):}
  \[
    \Delta T = -\frac{an}{C_{V,m}}\left(\frac{1}{V_1} - \frac{1}{V_2}\right) = -\frac{(0.3658)(10.00)}{28.94}\left(\frac{1}{5.000\times 10^{-3}} - \frac{1}{5.000\times 10^{-2}}\right) = -(0.12640)(200.0 - 20.0) = \mathbf{\SI{-22.75}{\kelvin}}
  \]
  Thus the $\mathrm{CO_2}$ cools to $T_2 = 350.0 - 22.75 = \mathbf{\SI{327.25}{\kelvin}}$ (whereas an ideal gas would have $\Delta T = 0$).
\end{enumerate}
\verifiedby{Total internal energy conservation: $\Delta U_T = nC_{V,m}\Delta T = \SI{-6584.4}{\joule}$ balances potential energy gain $\Delta U_V = an^2(1/V_1 - 1/V_2) = +\SI{6584.4}{\joule}$ (\texttt{check\_sec06\_09.py::ex06\_6}).}
\end{workedexample}

% ============================================================================
\section{Multi-Compartment Piston-Cylinder Equilibrium \& Coupled Gas Mechanics}\label{sec:06.7}
% ============================================================================
\idxtopic{piston-cylinder equilibrium}\idxtopic{multi-compartment cylinder}\idxtopic{spring-loaded piston}

Multi-compartment piston problems (\texttt{2018 MIT E3 Q24}, \texttt{2019 GGSO E4}) are premier tiebreakers at invitationals and Nationals \src{S25,S27,S28}. Solving them without calculus requires matching the mechanical and thermal constraints of the dividing piston.

\begin{formulacard}[Master Ledger for Coupled Piston-Cylinder Systems \src{S13,S27,S33}]
\begin{enumerate}
  \item \textbf{Case A — Free, Frictionless, \emph{Thermal-Conducting} (Diathermal) Piston Inside an Insulated Rigid Cylinder of Total Volume $V_{\mathrm{tot}} = V_1 + V_2$ (\texttt{2018 MIT E3 Q24}):}
  \begin{itemize}
    \item \emph{Equilibrium conditions:} Thermal equilibrium enforces $T_{1f} = T_{2f} = T_f$; mechanical equilibrium enforces $P_{1f} = P_{2f} = P_f$.
    \item \emph{Total internal energy conservation ($Q_{\mathrm{tot}}=0, W_{\mathrm{tot}}=0$):}
    \begin{equation}\label{eq:06:piston-diathermal-T}
      \Delta U_1 + \Delta U_2 = n_1 C_{V,m,1}(T_f - T_{1i}) + n_2 C_{V,m,2}(T_f - T_{2i}) = 0 \implies \boxed{T_f = \frac{n_1 C_{V,m,1} T_{1i} + n_2 C_{V,m,2} T_{2i}}{n_1 C_{V,m,1} + n_2 C_{V,m,2}}}
    \end{equation}
    \idxformula{Two-compartment diathermal piston final temperature: $T_f = (n_1C_{V,1}T_{1i}+n_2C_{V,2}T_{2i})/(n_1C_{V,1}+n_2C_{V,2})$}
    \item \emph{Final volumes and pressure:} Since $P_f = n_1 R T_f / V_{1f} = n_2 R T_f / V_{2f}$, the volume ratio equals the mole ratio:
    \begin{equation}\label{eq:06:piston-diathermal-V}
      \boxed{V_{1f} = \frac{n_1}{n_1 + n_2}V_{\mathrm{tot}}, \qquad V_{2f} = \frac{n_2}{n_1 + n_2}V_{\mathrm{tot}}, \qquad P_f = \frac{(n_1 + n_2)R T_f}{V_{\mathrm{tot}}}}
    \end{equation}
    \item \emph{Equal-$\gamma$ Constant-Pressure Theorem:} Using $n_k C_{V,m,k} T_{ki} = \frac{P_{ki}V_{ki}}{\gamma_k - 1}$, if both gases have the \textbf{same adiabatic index} ($\gamma_1 = \gamma_2 = \gamma$) and start in mechanical equilibrium ($P_{1i} = P_{2i} = P_i$), then:
    \begin{equation}\label{eq:06:piston-const-P}
      U_{\mathrm{tot}} = \frac{P_{1i}V_{1i} + P_{2i}V_{2i}}{\gamma - 1} = \frac{P_i V_{\mathrm{tot}}}{\gamma - 1} = \frac{P_f V_{\mathrm{tot}}}{\gamma - 1} \implies \boxed{P_f = P_i \quad (\text{pressure never changes!})}
    \end{equation}
  \end{itemize}
  \item \textbf{Case B — Free, Frictionless, \emph{Adiabatic} (Insulating) Piston with Heat $Q_{\mathrm{in}}$ Added to Side 1:}
  Side~2 undergoes a \textbf{reversible adiabatic compression} ($Q_2 = 0$), while mechanical equilibrium enforces $P_{1f} = P_{2f} = P_f$:
  \begin{equation}\label{eq:06:piston-adiabatic}
    P_{2i}V_{2i}^{\gamma_2} = P_f V_{2f}^{\gamma_2}, \qquad T_{2f} = T_{2i}\left(\frac{V_{2i}}{V_{2f}}\right)^{\gamma_2 - 1}, \qquad V_{1f} = V_{\mathrm{tot}} - V_{2f}, \qquad Q_{\mathrm{in}} = \Delta U_1 + \Delta U_2 = \frac{P_f V_{1f} - P_{1i}V_{1i}}{\gamma_1 - 1} + \frac{P_f V_{2f} - P_{2i}V_{2i}}{\gamma_2 - 1}
  \end{equation}
  \item \textbf{Case C — Vertical Heavy Piston (Mass $m_p$, Area $A$) and Linear Spring (Stiffness $k_{\mathrm{sp}}$):}
  \begin{equation}\label{eq:06:piston-spring}
    P_{\mathrm{gas}}(V) = P_{\mathrm{atm}} + \frac{m_p g_n}{A} + \frac{k_{\mathrm{sp}}(V - V_0)}{A^2} \quad (\text{straight line on } P\text{--}V \text{ plot} \implies W_{\mathrm{by}} = \frac{P_1 + P_2}{2}(V_2 - V_1))
  \end{equation}
  \idxformula{Spring-loaded piston pressure and work: $P(V) = P_0 + k_{\mathrm{sp}}(V-V_0)/A^2$, $W_{\mathrm{by}} = \frac{1}{2}(P_1+P_2)\Delta V$}
\end{enumerate}
\end{formulacard}

\begin{workedexample}[Two-Compartment Cylinder with Conducting Piston: Helium vs.\ Nitrogen]{2018 MIT E3 Q24 Archetype \textbar\ Tier 3}{ex:06.7}
A rigid, perfectly insulated horizontal cylinder of total internal volume $V_{\mathrm{tot}} = \SI{30.00}{\liter}$ is divided into two compartments by a thin, frictionless, thermally conducting (diathermal) piston. Initially, the piston is pinned in place: Compartment~1 ($V_{1i} = \SI{10.00}{\liter}$) holds $n_1 = \SI{2.000}{\mole}$ of monatomic Helium ($\gamma_1 = 5/3$, $C_{V,m,1} = \frac{3}{2}R$) at $T_{1i} = \SI{600.0}{\kelvin}$, while Compartment~2 ($V_{2i} = \SI{20.00}{\liter}$) holds $n_2 = \SI{3.000}{\mole}$ of diatomic Nitrogen ($\gamma_2 = 7/5$, $C_{V,m,2} = \frac{5}{2}R$) at $T_{2i} = \SI{300.0}{\kelvin}$. The pin is released and the system reaches full mechanical and thermal equilibrium. Find (a) the initial pressures $P_{1i}$ and $P_{2i}$, (b) the final temperature $T_f$, (c) the final compartment volumes $V_{1f}, V_{2f}$ and final pressure $P_f$, and (d) the work done and heat transferred between the two compartments.
\par\smallskip\noindent
\textbf{Solution.}
\begin{enumerate}
  \item \textbf{(a) Initial compartment pressures ($R = \SI{8.314463}{\liter\kilo\pascal\per\mole\per\kelvin}$):}
  \[
    P_{1i} = \frac{n_1 R T_{1i}}{V_{1i}} = \frac{(2.000)(8.314463)(600.0)}{10.00} = \mathbf{\SI{997.74}{\kilo\pascal}}, \qquad P_{2i} = \frac{n_2 R T_{2i}}{V_{2i}} = \frac{(3.000)(8.314463)(300.0)}{20.00} = \mathbf{\SI{374.15}{\kilo\pascal}}
  \]
  \item \textbf{(b) Final equilibrium temperature $T_f$ from \eref{eq:06:piston-diathermal-T}:}
  With $n_1 C_{V,m,1} = 2.000(1.5R) = 3.000\,R$ and $n_2 C_{V,m,2} = 3.000(2.5R) = 7.500\,R$:
  \[
    T_f = \frac{(3.000\,R)(600.0) + (7.500\,R)(300.0)}{3.000\,R + 7.500\,R} = \frac{1800.0 + 2250.0}{10.500} = \frac{4050.0}{10.500} = \mathbf{\SI{385.714}{\kelvin}}
  \]
  \item \textbf{(c) Final volumes and final pressure $P_f$ from \eref{eq:06:piston-diathermal-V}:}
  \[
    V_{1f} = \frac{2.000}{2.000 + 3.000}(\SI{30.00}{\liter}) = \mathbf{\SI{12.00}{\liter}}, \qquad V_{2f} = \frac{3.000}{5.000}(\SI{30.00}{\liter}) = \mathbf{\SI{18.00}{\liter}}
  \]
  \[
    P_f = \frac{(n_1 + n_2)R T_f}{V_{\mathrm{tot}}} = \frac{(5.000)(8.314463)(385.714)}{30.00} = \mathbf{\SI{534.50}{\kilo\pascal}}
  \]
  \item \textbf{(d) Internal energy transfer ($\Delta U_1 = -\Delta U_2$):}
  \[
    \Delta U_1 = n_1 C_{V,m,1}(T_f - T_{1i}) = (3.000)(8.314463)(385.714 - 600.0) = \mathbf{\SI{-5.3450}{\kilo\joule}}, \qquad \Delta U_2 = +\mathbf{\SI{5.3450}{\kilo\joule}}
  \]
  Helium transfers $\SI{5.345}{\kilo\joule}$ of total energy (combined expansion work + conducted heat) to the Nitrogen compartment.
\end{enumerate}
\verifiedby{Pressure balance check: $P_{1f} = n_1 R T_f / V_{1f} = 2(8.314463)(385.714)/12.00 = \SI{534.50}{\kilo\pascal} = P_{2f}$ (\texttt{check\_sec06\_09.py::ex06\_7}).}
\end{workedexample}

\begin{trapbox}[Top 6 Gas-Law \& Kinetic-Theory Exam Traps (\S6 Diagnostic Matrix)]
\par\smallskip\noindent
\begin{errtab}
Using $M = \SI{28.0}{\gram\per\mole}$ in $v_{\mathrm{rms}} = \sqrt{3RT/M}$ & Off by a factor of $\sqrt{1000} = 31.62$ (gives $\SI{16.3}{\meter\per\second}$ instead of $\SI{517}{\meter\per\second}$!).\idxerr{grams vs. kilograms in molar mass $M$ for $v_{\mathrm{rms}}$} & Always convert $M$ to $\si{\kilogram\per\mole}$ ($\mathrm{N_2}: \SI{0.02801}{\kilogram\per\mole}$); room-$T$ gas speeds are $\SIrange{250}{1900}{\meter\per\second}$. \\
Substituting $P_{\mathrm{gauge}}$ into $PV = nRT$ & Underestimates $P_{\mathrm{abs}}$ by $\SI{101.325}{\kilo\pascal}$, corrupting $n, V, T$ by $50\text{--}100\%$. & Add $P_{\mathrm{atm}} = \SI{101.325}{\kilo\pascal}$ ($\SI{14.7}{\text{psi}}$) whenever the problem states ``gauge'' or ``manometer height.'' \\
Forgetting $P_{\mathrm{sat,H_2O}}(T)$ in eudiometry & Overestimates dry gas moles by $2\text{--}5\%$ at room $T$ (and $>10\%$ above $\SI{45}{\celsius}$). & Subtract $P_{\mathrm{sat,H_2O}}(T)$ (\tref{tab:06:vdw-constants} / App.~A) and any liquid-column $\rho g h$ before using $PV=nRT$. \\
Averaging $\gamma_{\mathrm{mix}} = \sum x_i \gamma_i$ for a mixture & Violates energy additivity because $\gamma = C_P/C_V$ is a nonlinear ratio. & Compute $C_{V,m,\mathrm{mix}} = \sum x_i C_{V,m,i}$ and $C_{P,m,\mathrm{mix}} = C_{V,m,\mathrm{mix}} + R$ first, then take the ratio. \\
Confusing $K_{\mathrm{trans,avg}} = \frac{3}{2}k_B T$ with $U = \frac{f}{2}nRT$ & Gives $\frac{5}{2}k_B T$ for the \emph{translational} kinetic energy of $\mathrm{N_2}$ instead of $\frac{3}{2}k_B T$.\idxerr{translational kinetic energy vs. total internal energy for diatomic gas} & Translational kinetic energy is \textbf{always} $\frac{3}{2}k_B T$ per molecule regardless of whether the gas is monatomic or polyatomic! \\
Inverting Graham's law for \emph{mass} effusion rate & Predicts $\dot{m}_{\mathrm{He}} > \dot{m}_{\mathrm{O_2}}$ instead of $\dot{m}_1/\dot{m}_2 = \sqrt{M_1/M_2}$. & Molar rate $\dot{n} \propto 1/\sqrt{M}$, so mass rate $\dot{m} = M\dot{n} \propto \sqrt{M}$ (heavier gas has higher mass flux!). \\
\end{errtab}
\end{trapbox}

\begin{devicebox}[2027 Collector Air-Pocket Thermal Expansion, Buoyancy \& Dead-Air Mean Free Path]
\idxdev{air-gap convective onset and gas thermal conductivity}
In the 2027 Heat-Lamp Collector (\sref{sec:15}), trapped air pockets heated from $\SI{20}{\celsius}$ ($\SI{293.15}{\kelvin}$) to $\SI{65}{\celsius}$ ($\SI{338.15}{\kelvin}$) possess volumetric expansion coefficient $\beta_{V,\mathrm{air}} = 1/T \approx \SI{3.17e-3}{\per\kelvin}$, decreasing air density by $\Delta\rho/\rho \approx -13.3\%$ (from $\SI{1.204}{\kilogram\per\cubic\meter}$ to $\SI{1.044}{\kilogram\per\cubic\meter}$). If an insulation cavity is sealed hermetically, constant-volume heating raises internal pressure by $\Delta P = P_{\mathrm{atm}}(\Delta T/T_1) \approx \SI{+15.6}{\kilo\pascal}$ ($\SI{2.26}{\text{psi}}$), which can bow or pop off an unvented BoPET glazing lid (\texttt{Clarification \#5})!\idxdev{sealed glazing pressure bulging from Charles/Amontons law} Always leave a pinhole pressure-equalization vent in structural air baffles while keeping dead-air layer thickness $\delta \le \SI{11}{\milli\meter}$ to suppress Rayleigh--B\'enard buoyancy plumes (\sref{sec:11}).
\end{devicebox}

\alsosee{Quick Sheet \texttt{QS-05} (\sref{sec:01.5}); First Law \& Polytropic Processes (\sref{sec:07}); Ideal-Gas Entropy \& Mixing (\sref{sec:08}); Gas Power Cycles (\sref{sec:09}); Dead-Air Convection (\sref{sec:11}); Master Gas Tables (\sref{app:A}).}
